\documentclass[11pt]{article}
\usepackage{series}
\usepackage{booktabs}
\usepackage{array}
\usepackage{tabularx}

\newcommand{\cB}{\mathcal{B}}
\newcommand{\cC}{\mathcal{C}}
\newcommand{\cE}{\mathcal{E}}
\newcommand{\cU}{\mathcal{U}}
\newcommand{\cR}{\mathcal{R}}
\newcommand{\alphap}{\alpha'}
\newcolumntype{L}[1]{>{\raggedright\arraybackslash}p{#1}}
\renewenvironment{abstract}
  {\begin{center}\textbf{Abstract}\end{center}\begin{quote}\small}
  {\end{quote}}

\title{\textbf{Worldsheet and Spacetime Consistency as a Conditional Diagnostic Square:}\\
\large Sigma-model beta functions, target-space equations, and the MTT source contract}
\author{Peter Nero}
\date{July 2026, Version 2}

\begin{document}
\maketitle

\begin{abstract}
Perturbative string theory supplies a precise relation between two-dimensional
Weyl consistency and target-space field equations.  At leading order the
metric beta function contains the Ricci tensor, but the complete statement
also involves the antisymmetric tensor, dilaton, renormalization scheme,
higher powers of alpha-prime, string loops, and the domain of the worldsheet
theory.  It is therefore inaccurate to identify General Relativity and string
theory as theories, or to infer their equivalence merely because their
consistency equations meet in a low-energy corner.  This paper gives the
correct technical bridge and states what Modal Triplet Theory (MTT) would
have to add.  We define a worldsheet diagnostic, a target-space
Euler--Lagrange diagnostic, two source encodings, and an explicit comparison
operator.  A controlled diagnostic-square theorem proves exact and
approximate transfer of solutions when these maps intertwine and the
comparison operator has the required inverse estimate.  A counterexample
shows that coincident fixed points alone do not establish this bridge.
The standard sigma-model result supplies a perturbative instance after field
redefinitions; it does not supply an upstream MTT source.  We finish with a
twelve-row same-source contract and a separate error ledger.  Current q79
work provides five rows, partially provides two, and leaves five open.
Accordingly, the paper establishes a rigorous conditional comparison
framework, not a derivation of string theory, pure Einstein gravity, or their
ontological identity from MTT.
\end{abstract}

\section*{Version 2 Revision Note}
\begin{description}[leftmargin=2.8cm,style=nextline]
\item[Supersedes] Version 1.0, released in January 2026.
\item[Reason] The previous version identified GR and string theory as the
same constraint, assumed that an MTT proper-time flow was conjugate to
worldsheet RG, identified \(\alpha'\) with an inverse spectral gap, and
treated common fixed points as a proof of equivalence.
\item[Resolution] Version 2 separates the established perturbative
worldsheet-to-target result from the proposed MTT common-source explanation.
It proves a typed conditional diagnostic-square theorem, gives a no-go
counterexample, and records every approximation and source obligation.
\item[Retained content] The central intuition survives in narrower form:
worldsheet and spacetime equations may be two controlled diagnostics of one
selected background record.
\item[Open boundary] MTT has not yet supplied one selected q79 physical
visible--hidden worldsheet theory, upper action, and comparison map that
closes the twelve-row contract.
\end{description}

\tableofcontents

\section{What the familiar statement actually means}

The phrase ``GR falls out of string theory'' compresses several distinct
claims.  Perturbative string theory starts with a two-dimensional quantum
field theory whose couplings are target-space fields.  Quantum Weyl
invariance constrains those couplings.  In a perturbative scheme, the
constraints agree with equations derived from a target-space effective
action.  In a further low-energy and field-content reduction, part of that
action resembles Einstein gravity.

This chain is important, but it does not say that the theories are identical.
Their variables, observables, dimensions, quantization data, and regimes are
different.  The careful statement is:
\[
\begin{array}{c}
\text{worldsheet Weyl-anomaly coefficients: }\bar\beta^i=0\\[0.35em]
\Updownarrow\quad\text{at a declared order, scheme, and field content}\quad\\[0.35em]
\text{target-space equations: }\delta S_{\mathrm{eff}}/\delta\varphi^i=0 .
\end{array}
\]
The arrow concerns diagnostics on a shared background-coupling space.  It is
not an identity between complete theories.

\subsection{Three levels that must not be conflated}

\begin{description}[leftmargin=3.2cm,style=nextline]
\item[Standard bridge] A perturbative equivalence between Weyl-anomaly
coefficients and target effective equations, modulo field redefinitions and
at a stated order.
\item[MTT encoding] A proof that both diagnostics are images of one declared
upper source under explicit maps.
\item[MTT selection] A proof that the upper source, branch, action,
normalizations, and physical state are selected without importing the desired
lower solution.
\end{description}

This paper formalizes the second level conditionally.  It neither reproves all
of the first nor claims the third.

\section{The standard worldsheet-to-target bridge}

\subsection{Worldsheet data}

For a bosonic background record
\[
\varphi=(g_{\mu\nu},B_{\mu\nu},\Phi,\ldots)
\]
on a target \(Y\), the Euclidean sigma-model action contains
\begin{align}
S_\Sigma[X;\varphi]
={}&\frac{1}{4\pi\alphap}\int_\Sigma
\sqrt{h}\,h^{ab}g_{\mu\nu}(X)
\partial_aX^\mu\partial_bX^\nu\,d^2\sigma \nonumber\\
&+\frac{i}{4\pi\alphap}\int_\Sigma
\epsilon^{ab}B_{\mu\nu}(X)
\partial_aX^\mu\partial_bX^\nu\,d^2\sigma
+\frac{1}{4\pi}\int_\Sigma\sqrt h\,
R^{(2)}\Phi(X)\,d^2\sigma ,
\label{eq:sigma-action}
\end{align}
with convention-dependent normalizations.  A complete superstring or
heterotic model also needs worldsheet fermions, ghosts, spin structures,
GSO data, gauge-bundle couplings, a measure, and global consistency.

At leading order in one common convention, the metric anomaly coefficient has
the form
\begin{equation}
\bar\beta^g_{\mu\nu}
=\alphap\left(
R_{\mu\nu}+2\nabla_\mu\nabla_\nu\Phi
-\frac14H_{\mu\rho\sigma}H_\nu{}^{\rho\sigma}
\right)+O(\alphap^2),
\label{eq:metric-beta}
\end{equation}
where \(H=dB\) is modified in heterotic theory by the appropriate
Chern--Simons terms.  The \(B\)-field and dilaton have their own anomaly
coefficients.  Thus \(\bar\beta^g=0\) is not generally the vacuum Einstein
equation.

\subsection{Target-space data}

At string tree level and leading derivative order, the NS--NS part of a
target effective action is schematically
\begin{equation}
S_{\mathrm{eff}}
=\frac{1}{2\kappa^2}\int_Y
\sqrt{-g}\,e^{-2\Phi}
\left[
R+4|\nabla\Phi|^2-\frac1{12}|H|^2+O(\alphap)
\right]d^Dx .
\label{eq:effective-action}
\end{equation}
Gauge, fermion, source, compactification, and loop terms must be added for
the intended string.  After compatible field redefinitions, the
Euler--Lagrange equations of \eqref{eq:effective-action} agree with the
worldsheet Weyl conditions at the calculated order
\cite{CallanEtAl1985,MetsaevTseytlin1987,HullTownsend1986}.

Vacuum Einstein gravity appears only after further restrictions, for example
constant dilaton, vanishing flux and sources, suitable dimension reduction,
and neglect of higher-derivative and loop effects.  These are hypotheses, not
consequences of writing down \eqref{eq:sigma-action}.

\subsection{Scheme and field-redefinition dependence}

Beta functions are coordinates on a space of couplings.  Local field
redefinitions change their components and change the representative
effective action while preserving the appropriate on-shell physics.
Accordingly, comparison must either fix a common scheme or include the
field-redefinition map.  Equality of two unlabelled formulae is not a
scheme-independent theorem.

\section{A typed diagnostic square}

\subsection{Objects and maps}

Let \(\cU\) be a declared upper-source space.  Let
\(\cC_{\mathrm{ws}}\) be the space of complete worldsheet coupling records
and \(\cC_{\mathrm{st}}\) the space of target effective-field records.
Encoding maps
\[
E_{\mathrm{ws}}:\cU\to\cC_{\mathrm{ws}},
\qquad
E_{\mathrm{st}}:\cU\to\cC_{\mathrm{st}}
\]
must be independently defined.  They are not names for fitting the same
lower data twice.

Let
\[
D_{\mathrm{ws}}:\cC_{\mathrm{ws}}\to\cB_{\mathrm{ws}},
\qquad
D_{\mathrm{st}}:\cC_{\mathrm{st}}\to\cB_{\mathrm{st}}
\]
denote the full Weyl-anomaly and target Euler--Lagrange diagnostics.
The spaces \(\cB_{\mathrm{ws}}\) and \(\cB_{\mathrm{st}}\) carry declared
norms on a domain \(\cU_0\subseteq\cU\).  A comparison family
\[
J_u:\cB_{\mathrm{st}}\to\cB_{\mathrm{ws}}
\]
includes conventions, field redefinitions, gauge quotients, and any
dimensional reduction used in the comparison.

\begin{definition}[Controlled diagnostic square]
The data above form a controlled diagnostic square on \(\cU_0\) when
\begin{equation}
D_{\mathrm{ws}}(E_{\mathrm{ws}}u)
=J_uD_{\mathrm{st}}(E_{\mathrm{st}}u)+r(u),
\qquad
\|r(u)\|_{\mathrm{ws}}\leq\varepsilon(u),
\label{eq:square}
\end{equation}
and \(J_u\) is injective on the compared diagnostic subspace with
\[
\|v\|_{\mathrm{st}}\leq M(u)\|J_uv\|_{\mathrm{ws}}.
\]
The square is exact when \(r=0\).
\end{definition}

The inverse estimate matters.  Without it, a small worldsheet residual can
hide a large target residual in a poorly conditioned or discarded direction.

\begin{theorem}[Controlled worldsheet--spacetime diagnostic transfer]
\label{thm:square}
Let the controlled diagnostic square hold at \(u\in\cU_0\).  Then
\begin{equation}
\|D_{\mathrm{st}}(E_{\mathrm{st}}u)\|_{\mathrm{st}}
\leq M(u)\left(
\|D_{\mathrm{ws}}(E_{\mathrm{ws}}u)\|_{\mathrm{ws}}
+\varepsilon(u)\right).
\label{eq:target-bound}
\end{equation}
If the square is exact and \(J_u\) is injective on the compared subspace,
worldsheet consistency implies the compared target equations.  If \(J_u\)
is also surjective and boundedly invertible on the full diagnostic spaces,
the two zero conditions are equivalent.  None of these conclusions identifies
the complete theories.
\end{theorem}

\begin{proof}
Rearranging \eqref{eq:square} gives
\[
J_uD_{\mathrm{st}}(E_{\mathrm{st}}u)
=D_{\mathrm{ws}}(E_{\mathrm{ws}}u)-r(u).
\]
Apply the inverse estimate and the triangle inequality to obtain
\eqref{eq:target-bound}.  For \(r=0\), injectivity transfers a worldsheet
zero to a target zero.  Surjectivity and bounded invertibility give the
reverse implication on the declared spaces.  The theorem compares diagnostic
values only, so it makes no statement about equality of states, observables,
or quantizations.
\end{proof}

\begin{corollary}[Perturbative string instance]
Suppose a fixed string model, renormalization scheme, field-redefinition
map, and perturbative order \(N\) give
\[
D_{\mathrm{ws}}=J D_{\mathrm{st}}+O(\alphap^{N+1})
\]
on a controlled background family.  Then Theorem~\ref{thm:square} transfers
the residual with \(\varepsilon=O(\alphap^{N+1})\), together with whatever
string-loop, compactification, and analytic errors are separately present.
\end{corollary}

This corollary organizes the standard result.  It does not prove that an MTT
source emitted the string model.

\section{Why matching fixed points is insufficient}

\begin{proposition}[Common-zero no-go]
\label{prop:no-go}
Equality of the zero sets of two diagnostics does not imply a controlled
diagnostic square with a uniformly invertible linear comparison.
\end{proposition}

\begin{proof}
On \(\mathbb{R}\), let \(D_1(x)=x\) and \(D_2(x)=x^3\).  Both have zero set
\(\{0\}\).  Any pointwise scalar comparison satisfying
\(D_2(x)=J_xD_1(x)\) for \(x\neq0\) has \(J_x=x^2\).  Its inverse norm
diverges as \(x\to0\).  Thus common fixed points do not provide the stability
or residual transfer required by Theorem~\ref{thm:square}.
\end{proof}

This elementary example exposes the defect in an argument of the form
``both flows have the same fixed point, therefore they encode the same
constraint.''  A second circular argument is to assume a conjugacy
\(F_{\mathrm{ws}}=U^{-1}F_{\mathrm{st}}U\) and then advertise fixed-point
transfer as a derivation.  Conjugacy does transfer fixed points, but the
research task is to construct and verify \(U\).  Naming it does not do so.

\section{What MTT would add}

MTT proposes an explanation upstream of the standard bridge: both lower
diagnostics may descend from one selected source.  That proposal becomes
mathematical only after the maps in \eqref{eq:square} are built from the same
source, with provenance.

\subsection{The same-source contract}

\begin{center}
\small
\begin{tabularx}{\textwidth}{@{}L{0.07\textwidth}X L{0.27\textwidth}@{}}
\toprule
Row & Required object & Current q79 status\\
\midrule
W1 & Time-oriented q79 target branch & Available\\
W2 & Fu--Yau charge and Green--Schwarz/Bianchi sector & Available\\
W3 & Visible curvature-level Green--Schwarz cancellation & Available at
curvature tier\\
W4 & Critical heterotic central-charge balance & Available universally\\
W5 & q79 low-energy GR and quantum-EFT parity limit & Available at its
declared tier\\
W6 & Global differential gerbe and full visible-cycle consistency & Open;
finite representative only\\
\bottomrule
\end{tabularx}

\medskip
\begin{tabularx}{\textwidth}{@{}L{0.07\textwidth}X L{0.27\textwidth}@{}}
\toprule
Row & Required object & Current q79 status\\
\midrule
W7 & All-orders-in-\(\alphap\) q79 target background & Open; present
background is first order\\
W8 & Exact q79 heterotic \((0,2)\) SCFT or complete beta functions &
Partial; base GLSM and local/topological data exist, but the physical bundle
and IR SCFT do not\\
W9 & Modular-invariant q79 GSO partition function and factorization &
Partial; finite torsion data exist, while analytic characters and GSO remain
open\\
W10 & q79-specific string-field vertices and BV master action & Open\\
W11 & Tadpole, vacuum-shift, infrared, and soft-state completion & Open\\
W12 & All-genus convergence or a nonperturbative definition & Open\\
\bottomrule
\end{tabularx}
\end{center}

The count is five available, two partial, and five open.  In the live
research ledger, \(B.\mathrm{QG}.01\) asks for all twelve rows on the same
physical nonpullback visible--hidden bundle.  \(B.\mathrm{ACTION}.01\)
separately asks for the selected upper differential or action whose
automorphisms and descendants reproduce the lower structures.  Neither
blocker is closed by the theorem in this paper.

\subsection{Why a projector and a gap are not enough}

A coherent projector can select a subspace, and a spectral gap can control
leakage from it.  Neither object alone determines a target manifold,
worldsheet measure, BRST complex, anomaly coefficients, effective action,
comparison map, or renormalization scheme.  Likewise, proper-time
representations and worldsheet Wilsonian RG are not automatically the same
flow.  To identify them one needs a proved intertwiner with domain and error
control.

The formerly asserted relation
\[
\alphap\sim\lambda_*^{-1}
\]
is therefore not used as an identity.  \(\alphap\) is the string tension
scale in the sigma model; \(\lambda_*\) may be a spectral control scale in an
MTT reduction.  A selected source theorem could relate them, but dimensional
compatibility and matching of one asymptotic order would still not determine
all coefficients.

\section{The complete error ledger}

A useful bridge reports errors by origin rather than hiding them in one
symbol.

\begin{center}
\small
\begin{tabular}{@{}p{0.22\textwidth}p{0.32\textwidth}p{0.35\textwidth}@{}}
\toprule
Source & Typical control & What failure means\\
\midrule
\(\alphap\) expansion & curvature and derivative bounds in string units &
higher-derivative terms are not small\\
String loops & powers of \(g_s=e^\Phi\) and genus estimates & tree-level
target action is insufficient\\
Worldsheet analysis & anomaly, modular, GSO, factorization, IR SCFT &
no complete perturbative string background\\
Scheme/redefinition & explicit local map and Jacobian bounds & component
beta functions cannot be compared directly\\
Compactification & Kaluza--Klein scale, moduli, warping, sources & the
lower-dimensional GR truncation is uncontrolled\\
MTT projection & leakage, domain, spectral and truncation estimates &
upper-to-lower descent is uncontrolled\\
Diagnostic comparison & residual \(r\) and inverse bound \(M\) & lower
zeros or small residuals do not transfer\\
Quantum/observable match & state, BRST/BV, renormalization, uncertainty &
equation-level parity is not physical equivalence\\
\bottomrule
\end{tabular}
\end{center}

These errors need not scale together.  In particular, a small projection
error says nothing by itself about string loops or modular consistency.

\section{Interpretation and consequences}

\subsection{What is genuinely unified}

The standard string result already unifies a set of worldsheet and
target-space consistency equations within one perturbative construction.
The MTT proposal is stronger in a different direction: it seeks one
preprojection object from which the worldsheet and spacetime records descend.
If the twelve-row contract and Theorem~\ref{thm:square}'s hypotheses are
realized from one selected q79 source, then the two diagnostics would be
certified views of that source.  This would explain their agreement without
claiming that a two-dimensional QFT and a spacetime effective theory are the
same mathematical object.

\subsection{What is not yet unified}

The present framework does not derive:
\begin{itemize}
\item a complete physical q79 heterotic worldsheet;
\item an exact \((0,2)\) IR SCFT or modular-invariant GSO sum;
\item the nonperturbative string theory or all-genus completion;
\item a selected upper MTT action;
\item pure four-dimensional Einstein gravity without reduction hypotheses;
\item equality of worldsheet and spacetime observables; or
\item a unique physical vacuum or branch.
\end{itemize}

It also does not prove that all quantum gravity theories are shadows of one
constraint.  That may be a research program, but each proposed diagnostic
square needs its own maps and certificate.

\subsection{Falsifiability}

Once source maps are explicit, the bridge can fail in informative ways.
The worldsheet record may violate modular invariance; the target equations
may contain a component outside the range controlled by \(J\); the inverse
bound may diverge; the \(\alphap\) or loop expansion may cease to be small;
or two purported descendants may have incompatible source hashes.  Any of
these failures blocks promotion.  Agreement obtained only after adjusting
the source to the desired lower equations is reconstruction, not prediction.

\section{Research program}

The shortest rigorous path is:
\begin{enumerate}
\item construct the physical q79 visible and hidden holomorphic bundles in
one positive HYM chamber;
\item close the global differential Green--Schwarz and gerbe data;
\item construct the heterotic \((0,2)\) worldsheet theory, analytic
characters, GSO projection, and factorization;
\item derive its anomaly coefficients in a declared scheme;
\item derive the target effective equations from the same source and at the
same order;
\item emit \(J\), \(r\), \(\varepsilon\), and the inverse estimate \(M\);
\item compare states and observables only after the equation-level square is
closed; and
\item keep the all-genus or nonperturbative completion as a separate tier.
\end{enumerate}

This order prevents a familiar loop: proving a formal transfer theorem again
while the physical source rows remain absent.  The theorem is now owned here.
Future progress must fill a row or improve a bound.

\section{Conclusion}

General Relativity and perturbative string theory are not the same theory.
The accurate mathematical relation is a controlled correspondence between
worldsheet Weyl diagnostics and target-space effective equations on a shared
background record, at a declared perturbative order and scheme.  Pure
Einstein gravity is a further low-energy corner.

MTT offers a potentially deeper explanation only if one selected upper source
emits both records and the comparison square.  Theorem~\ref{thm:square}
states exactly what that would prove and Proposition~\ref{prop:no-go} shows
why common fixed points do not suffice.  With five of twelve q79 worldsheet
rows available, two partial, and five open, the bridge is a serious,
well-typed program rather than a completed equivalence.

\begin{thebibliography}{99}

\bibitem{Friedan1980}
D.~Friedan,
``Nonlinear models in \(2+\epsilon\) dimensions,''
\emph{Physical Review Letters} \textbf{45} (1980), 1057--1060.
\url{https://doi.org/10.1103/PhysRevLett.45.1057}

\bibitem{CallanEtAl1985}
C.~G. Callan, D.~Friedan, E.~J. Martinec, and M.~J. Perry,
``Strings in background fields,''
\emph{Nuclear Physics B} \textbf{262} (1985), 593--609.
\url{https://doi.org/10.1016/0550-3213(85)90506-1}

\bibitem{HullTownsend1986}
C.~M. Hull and P.~K. Townsend,
``String effective actions from sigma-model conformal anomalies,''
\emph{Nuclear Physics B} \textbf{301} (1988), 197--223.
\url{https://doi.org/10.1016/0550-3213(88)90692-5}

\bibitem{MetsaevTseytlin1987}
R.~R. Metsaev and A.~A. Tseytlin,
``Order \(\alphap\) (two-loop) equivalence of the string equations of
motion and the sigma-model Weyl invariance conditions,''
\emph{Nuclear Physics B} \textbf{293} (1987), 385--419.
\url{https://doi.org/10.1016/0550-3213(87)90077-0}

\bibitem{NeroProjectionString}
P.~Nero,
\emph{A Projection-First Reframing of String Theory: Conditional
Encodings, Worldsheet Gates, and the q79 Boundary},
Version 2, 2026.

\bibitem{NeroStrominger}
P.~Nero,
\emph{Modal Triplet Theory to the Hull--Strominger System: A Conditional
Fixed-Point Correspondence},
Version 2, 2026.

\end{thebibliography}
\end{document}
